\kbvcontinued ಸಂಧಿ-ಇವು ಸ್ವರ ಸಂಧಿಗಳು. ವ್ಯಂಜನ ಸಂಧಿಗಳಲ್ಲಿ ಮುಖ್ಯವಾದವುಗಳು\par

\vspace{2.5pt plus 3pt minus 0.5pt}\noindent\hspace*{1.2em}\kbvnum{(೧)}\hspace{0.55em}ಜಶ್ತ್ವ ಸಂಧಿ, (೨) ಶ್ಚುತ್ವ ಸಂಧಿ, ಮತ್ತು (೩) ಅನುನಾಸಿಕ ಸಂಧಿ. ಈ ಸಂಧಿಗಳಿಗೆ ಅನ್ವಯಿಸುವ ಸೂತ್ರಗಳನ್ನು ಉದಾಹರಣೆಗಳ ಸಹಿತ ತಿಳಿಯೋಣ.\par

\kbvsubtitle{ಸ್ವರ ಸಂಧಿಗಳು}

\vspace{4pt plus 4pt minus 1pt}\noindent{\kbvsecfont\fontsize{11pt}{14.5pt}\selectfont\color{kbvaccentdark} (೧)\hspace{0.45em}ಸವರ್ಣದೀರ್ಘ ಸಂಧಿ :}\hspace{0.5em}(ಅ, ಆ); (ಇ, ಈ); (ಉ, ಊ); (ಋ, ೠ); (ಎ, ಏ); (ಒ, ಓ)-ಇವು ಒಂದೇ ಜಾತಿಯ ಸ್ವರಗಳನ್ನುಳ್ಳ ಆರು ಜೊತೆ ಅಕ್ಷರಗಳು. ಇವುಗಳಿಗೆ ಸವರ್ಣ ಸ್ವರಗಳೆಂದು ಹೆಸರು. ಈ ಹಿನ್ನೆಲೆಯಲ್ಲಿ “ದೇವಾಸುರೆ” ಎಂಬ ಶಬ್ದದಲ್ಲಿನ ಸಂಧಿಕಾರ್ಯವನ್ನು ನೋಡೋಣ.\par

\begin{kbvderiv}
\vspace{4pt plus 3pt minus 1pt}
\begin{center}\begin{adjustbox}{max width=\textwidth}
{\renewcommand{\arraystretch}{1.22}\setlength{\tabcolsep}{4pt}
\begin{tabular}{l c l}
ದೇವಾಸುರೆ & {\color{kbvgrey}$=$} & ದೇವ + ಅಸುರ \\
 & {\color{kbvgrey}$=$} & ದೇ + ವ್‌ + (ಅ + ಅ) + ಸುರ \\
\end{tabular}}
\end{adjustbox}\end{center}
\vspace{4pt plus 3pt minus 1pt}
\end{kbvderiv}

\indent ಇಲ್ಲಿ “ಅಕಾರಕ್ಕೆ” ಅಕಾರ ಪರವಾಗಿದೆ ; ಇವು ಸವರ್ಣ ಸ್ವರಗಳು.\par

\indent  = ದೇ + ವ್‌ + (ಆ) + ಸುರ\par

\indent ಇಲ್ಲಿ ಅ, ಅಗಳ ಸ್ಥಾನದಲ್ಲಿ “ಆ” ಕಾರ ಆದೇಶವಾಗಿ ಬಂದಿದೆ.\par

\indent  = ದೇವಾಸುರ (ವ್‌ + ಆ = ವಾ)\par

\indent ಇಲ್ಲಿ ಎರಡು ಸವರ್ಣ ಸ್ವರಗಳಿಗೆ ಒಂದೇ ದೀರ್ಧ ಸ್ವರವು ಆದೇಶವಾಗಿ ಬಂದಿದೆ. ಇದನ್ನು ಸವರ್ಣದೀರ್ಧ ಸಂಧಿ ಎನ್ನುವರು.\par

\begin{kbvrule}
\noindent{\kbvsecfont ಸೂತ್ರ :}\hspace{0.5em}ಸವರ್ಣ ಸ್ವರಗಳು ಪರವಾದಾಗ (ಒಂದರ ಮುಂದೆ ಮತ್ತೊಂದು ಬಂದಾಗ) ಅವೆರಡರ ಸ್ಥಾನದಲ್ಲಿ ಅದೇ ಜಾತಿಯ ದೀರ್ಧ ಸ್ವರವು ಆದೇಶವಾಗಿ ಬರುವುದು. ಇದನ್ನು ಸವರ್ಣ ದೀರ್ಧ ಸಂಧಿಯೆನ್ನುವರು.
\end{kbvrule}

\vspace{3pt plus 2pt}\noindent\hspace*{1.2em}\kbvlabel{ಉದಾ}\par\vspace{1pt}

\begin{kbvderiv}
\vspace{4pt plus 3pt minus 1pt}
\begin{center}\begin{adjustbox}{max width=\textwidth}
{\renewcommand{\arraystretch}{1.22}\setlength{\tabcolsep}{4pt}
\begin{tabular}{l c l}
ಗುರು + ಉಪದೇಶ & {\color{kbvgrey}$=$} & ಗುರೂಪದೇಶ .... (ಉ + ಉ = ಊ) \\
ಸ್ವರ + ಅಕ್ಷರ & {\color{kbvgrey}$=$} & ಸ್ವರಾಕ್ಷರ .... (ಅ + ಅ = ಆ) \\
ಗಿರಿ + ಈಶ & {\color{kbvgrey}$=$} & ಗಿರೀಶ .... (ಇ + ಈ = ಈ) \\
\end{tabular}}
\end{adjustbox}\end{center}
\vspace{4pt plus 3pt minus 1pt}
\end{kbvderiv}
